\subsection{Cronograma de ejecución de la mejora}

La ejecución de la propuesta técnica se organiza en seis fases desarrolladas entre octubre y diciembre de 2026. 
La Tabla~\ref{tab:cronograma_ejecucion} detalla la programación temporal de cada actividad.

\par
\noindent
\begin{minipage}{\linewidth}
    \tituloTablaAPA{tab:cronograma_ejecucion}{Cronograma de ejecución de la propuesta de mejora}

    \begingroup
    \renewcommand{\baselinestretch}{1}
    \footnotesize
    \setlength{\tabcolsep}{4pt}
    \renewcommand{\arraystretch}{1.3}

    \begin{tabularx}{\linewidth}{@{}p{3.8cm}Xcccccc@{}}
        \toprule
        \textbf{Fase / Etapa} & \textbf{Actividades principales} & \textbf{S1} & \textbf{S2} & \textbf{S3} & \textbf{S4} & \textbf{S5} & \textbf{S6} \\ 
        & & \footnotesize(Oct) & \footnotesize(Nov 1) & \footnotesize(Nov 2) & \footnotesize(Nov 3) & \footnotesize(Nov 4) & \footnotesize(Dic) \\ \midrule
        Fase 1: Análisis y diseño & Diagnóstico, levantamiento de requerimientos y diseño de base de datos. & X & & & & & \\ \addlinespace
        Fase 2: Sprint 1 & Autenticación RBAC, contratos compartidos e importación de empleados. & & X & & & & \\ \addlinespace
        Fase 3: Sprint 2 & Marcación web/CSV/Excel, gestión de asistencias y auditoría. & & & X & & & \\ \addlinespace
        Fase 4: Sprint 3 & Mantenimiento de cargos, consolidación de incidencias y solicitudes. & & & & X & & \\ \addlinespace
        Fase 5: Sprint 4 & Dashboard contable, consulta de empleados y novedades de nómina. & & & & & X & \\ \addlinespace
        Fase 6: Integración y cierre & Flujos n8n (boletas y dispersión), pruebas finales y despliegue. & & & & & & X \\ \bottomrule
    \end{tabularx}
    \endgroup

    \notaTabla{S1 a S6 corresponden a las semanas de trabajo programadas. Elaboración propia.}
\end{minipage}
\finBloqueAPA